\chapter{两个宇宙对照}

本附录以表格形式，汇总正交宇宙（黎曼型时空）与我们自身宇宙（洛伦兹型时空）在物理各层面的
系统性差异。表中每一行都可在正文相应章节找到详细讨论。为便于对照，我们把“洛伦兹宇宙”列
置于左，“黎曼宇宙”列置于右。

\begin{longtable}{p{2.9cm} p{5.7cm} p{5.9cm}}
\toprule
\textbf{项目} & \textbf{洛伦兹宇宙（我们的）} & \textbf{黎曼宇宙（正交）} \\
\midrule
\endfirsthead
\toprule
\textbf{项目} & \textbf{洛伦兹宇宙（我们的）} & \textbf{黎曼宇宙（正交）} \\
\midrule
\endhead
\bottomrule
\endlastfoot

时空度规 &
$\dd s^{2}=-\dd t^{2}+\dd x^{2}+\dd y^{2}+\dd z^{2}$，号差 $(-,+,+,+)$ &
$\dd s^{2}=\dd t^{2}+\dd x^{2}+\dd y^{2}+\dd z^{2}$，号差 $(+,+,+,+)$ \\

坐标地位 &
时间方向特殊；光锥把类时与类空分开 &
四个维度完全平等；无光锥、无类光矢量 \\

光速 &
存在普适常量 $c$，一切观察者一致 &
\textbf{无普适光速}；相速依赖波长，$\omega^{2}+|\mathbf{k}|^{2}=\wms^{2}$ \\

速度上限 &
任何速度不超过 $c$ &
\textbf{无速度上限} \\

能量--动量关系 &
$E^{2}-p^{2}=m^{2}$ &
$E^{2}+p^{2}=m^{2}$；(B.1) 节 \\

光子静止质量 &
零 &
\textbf{非零}；由 Proca 型方程描述 \\

静电场（库仑势） &
$1/r$，无质量光子 &
随空间\emph{振荡}（质量项使其不指数衰减） \\

光的偏振态 &
两个（自旋 $\pm1$） &
\textbf{三个}（垂直于波矢的三维子空间） \\

发光过程 &
产生光需\emph{消耗}能量 &
\textbf{产生光释放能量}（放热） \\

高速远行的时间 &
旅行者时间膨胀，比家乡人更\emph{年轻} &
旅行者经历更\emph{长}的岁月（绕行使时间更长） \\

引力：Schwarzschild 解 &
有视界，$r=2M$ 处出现黑洞 &
径向距离有限，\textbf{无视界、无黑洞} \\

引力波 &
存在（四维时空中的横波） &
\textbf{不存在}（黎曼度规下无相应自由度） \\

宇宙学 &
大爆炸、膨胀宇宙为常见图景 &
\textbf{很可能没有大爆炸}；均匀各向同性解另有形态 \\

热力学：温度 &
温度无上界（上限趋于“无穷热”） &
存在\emph{负温度}，可“比无穷热更热” \\

黑体谱 &
普朗克谱，总功率按 $T^{4}$ 标度 &
谱被截止于 $\wms$；经典总能量按 $T$ 线性标度，经典/量子差异很小 \\

生命与光合作用 &
植物\emph{吸收}光能制造养分 &
植物向夜空\emph{发出}自己的光来制造养分 \\

星际航行 &
需携带燃料/推进剂 &
\textbf{发光引擎}可无燃料、无推进剂地持续加速 \\

量子力学 &
Dirac 方程、Klein--Gordon 方程（洛伦兹号差） &
黎曼 Dirac 方程；旋量由 $\mathrm{SU}(2)\times\mathrm{SU}(2)$ 双覆盖描述 \\

\end{longtable}

\section*{结语}

正交宇宙与我们的宇宙之别，归根到底只源于时空度规中时间项的一个符号。然而这一个符号的翻转，
却连锁地改写了光学、电磁学、热力学、引力与量子力学的面貌：光失去了普适速度、获得了静止质量；
发光由吸能变为释能；黑洞、引力波与大爆炸都不复存在；而“向黑夜发光”的植物与“无燃料的引擎”
则重塑了整个文明的生存方式。Egan 的《正交》三部曲，正是一次以如此单一的几何改动为起点、
把物理推论一路推到底的系统性思想实验。
